% ============================================================
%  CHAPTER III
%  Transmission of Visual Perception on a Plane
%  Source pp. 52–79
% ============================================================
\chapter{Transmission of Visual Perception on a Plane}
\markboth{Chapter III. Transmission of Visual Perception on a Plane}{Chapter III}

\srcpage{52}

In what way and to what degree is it in principle possible to reproduce the external
world exactly as it is visually perceived, considering only the geometry of the image
(setting aside shadow modelling and colour)?  In what way and to what degree is a
verbatim accurate transmission of the visible geometry of external space onto the
picture plane possible?

The initial requirement adopted in the present chapter~--- the aspiration to mathematical
exactness in the transmission of visible geometry~--- will be progressively relaxed in the
subsequent chapters, which will bring the examination of the problem of spatial
construction closer to the deeper sources of artistic creativity.

The construction of the appropriate system of perspective makes it possible to supply the
picture with a geometric foundation, analogously to the way in which the method of
orthogonal projections provides the geometric basis of a technical drawing.

The usual definitions of perspective~--- as the geometric discipline of conveying the
volumetric-spatial properties of objects on a plane~--- are constructed so as to contain
from the outset the assumption of central projection as the basis of perspective.  The
deficiency of such a definition is that it relies from the start on a method whose
expediency still remains to be demonstrated.  We shall therefore use another definition,
given in the perspective textbook of V.\,E.\,Peterson and of more general character:
\emph{perspective is the discipline of image methods corresponding to visual
perception.}\footnote{%
  V.\,E.~Peterson. \emph{Perspective}. Moscow, 1970, p.\,5.  The terms of the theory
  of linear perspective used below coincide with those given in this textbook.}
Although Peterson's textbook itself presents only the method of central projection, the
definition quoted permits a broader interpretation.

As already shown, visual perception is a complex process founded on the coordinated
work of the eye and the brain.  As a result of this work there arises the visible image of
the object contemplated.  Vision is therefore the product of the joint work of eye and
brain; a person is more capable of ``seeing without eyes'' (as in a dream) than ``without
a brain'' (as when unconscious).  In consequence, such concepts as \emph{to see},
\emph{vision}, \emph{visible image}, and the like will throughout this book denote only
the result of the joint work of eye and brain (and will never signify the retinal image
or other intermediate images in the process of visual perception).  Consequently,
perspective is a method of depicting on a plane the forms of objects, their mutual
arrangement, and so on, in a way that transmits to the viewer the visible image of the
external world.

\srcpage{53}

We shall call a \emph{system of scientific perspective} one that is obtained
mathematically solely from the objective laws of the psychology of human visual
perception.  Naturally, such a scientific system cannot and should not replace the
creativity of the artist.  An image obtained by its means will certainly have no artistic
value, although it may prove useful both to the artist and to the art historian for a deeper
understanding of the laws of perspective construction.  Within the framework of such a
system of perspective one may attempt to obtain a mathematically irreproachable image
of the external space on the picture plane.

The construction would be \emph{ideal} (in the mathematical sense, not in the artistic)
if, looking at the picture, the viewer were unable to distinguish the visual images that
arose from those arising from contemplating the corresponding objective external space.

Let us introduce the formal concept of a ``geometrically ideal image,'' corresponding to
the requirements formulated above.  If one speaks only of the quantitative correspondence
of geometric characteristics, this means that when in visual perception two segments
appear equal (regardless of their true sizes in external space), they should also appear
equal when the picture is contemplated; similarly, a segment that appears twice as long
should also appear twice as long in the picture, and so on.  This applies not only to
linear magnitudes but also to angles.  Lines that appear parallel, or perpendicular, or
inclined to one another at some angle~$\alpha$ in perceptual space should likewise appear
parallel, perpendicular, or intersecting at the same angle~$\alpha$ when the picture is
contemplated.

It is evident that such an ``ideal image'' presupposes a system of ``deceptions'' of
human visual perception, an imitation of the visual perception of space in the course of
contemplating the picture.  In accordance with the two-stage character of the visual
perception of external space, two fundamentally different routes open here, both leading
to the same goal.

1.~One may depict the retinal image on the canvas.\footnote{%
  Strictly speaking, the image constructed on the canvas must be such that, when
  viewed, an image on the retina is formed that is indistinguishable from the retinal
  image of the real space depicted on the canvas.  It is easy to see that such an image
  must be constructed by central projection, taking the eye as the centre of projection,
  i.e.\ (assuming a flat picture) by using the usual linear perspective.  It is important
  to note that in creating the picture one need not account for the actual shape of the
  retina (which is concave), since to produce the same effect from viewing real space
  and a picture one must obtain identical retinal images~--- and their actual shape is
  immaterial.  In this connection, various refinements to the system of linear perspective
  that take into account the curvature of the retina (see, e.g.: E.\,Panofsky.
  ``Die Perspektive als symbolische Form.'' \emph{Vortr\"age der Bibliothek Warburg},
  1924--1925; idem. \emph{Aufs\"atze zu Grundlagen der Kunstwissenschaft}.
  Berlin, 1964) are devoid of sense.  Wherever it is said below that the retinal image
  is depicted on the canvas, the conditions formulated in this note are always understood
  to hold.  For a more detailed account see Appendix~1.}
  Since a person never actually sees it, this can be done only by using the
corresponding optical-geometric regularities.  On such an image all depth cues must be
emphasized as strongly as possible in order to ``activate'' the size-constancy mechanism,
which must then subconsciously transform the retinal image arising from contemplating
the picture into the appropriate perceptual image~--- the one that would arise from
contemplating the space depicted in the picture~--- and to stimulate the action of the
form-constancy mechanism.

2.~One may depict on the canvas the visible image of the external world, i.e.\ that
which, being based on the retinal image, has already been transformed by the
regularities of visual perception.  In this case, however, one must strive to ensure that
when the picture is contemplated the constancy mechanisms are ``activated'' not ``at full
power'' but with due discretion and selectivity, since their action has already been
accounted for in the image, and a repeated geometric transformation of the retinal image
arising from contemplating the picture may distort the image that would arise from
contemplating the real space depicted.

\srcpage{54}

If one recalls the simplified scheme of the visual perception process given above, now
refined to: \emph{objective external space}~$\to$ \emph{retinal image}~$\to$
\emph{perceptual space}, the difference between the two proposed methods of
depicting perceptual space on a canvas becomes clear.  In the first case, the first stage
of visual perception is depicted on the canvas, and the transition to the formation of a
full visual image is substantially linked to the second arrow in the scheme~--- representing
the action of the constancy mechanisms~--- i.e.\ the perception of the picture must in a
certain sense also be two-staged.  In the second case, when perceptual space is depicted
directly, the transformations of the retinal image that alter its ``geometry'' may prove
undesirable~--- perception must in the same conditional sense be one-staged.

In such abstract reasoning both methods of conveying space are completely equivalent,
since they both ultimately lead to the emergence of the required image of perceptual
space.  The situation changes, however, if one considers the problem of the practical
realisation of both envisaged ideal schemes.

The first method, based on depicting the retinal image, depends essentially on the
effectiveness of processes equivalent to the action of the visual constancy mechanisms,
which must be activated by the picture.  This concerns in the first place the possibility
of reproducing on the canvas all the depth cues.  The monocular depth cues listed above
are quite reproducible by the artist on the canvas, which cannot be said of the binocular
ones.  Indeed, when a picture is contemplated, convergence will only indicate that all
points of the picture plane are equidistant from the viewer~--- but this can sharply
contradict the image if the artist has aimed to convey deep space.  It is easy to see that
disparity will lead to a completely analogous effect, and tactile-kinesthetic
representations will likewise contradict the assumption of whatever depth of space the
artist has shown.  Consequently, the desired complete transformation of the retinal image
arising from contemplating the picture into the required perceptual space is fundamentally
impossible.  Thus, the creation of a ``geometrically ideal image'' by the first of the
schemes is impossible.

The second method is free of these defects, since it reduces to a direct depiction of the
geometry of perceptual space.  In this case the binocular depth cues and
tactile-kinesthetic representations can even promote correct perception of the picture,
as they will effectively ``inhibit'' the subconscious geometric transformations of what is
depicted by processes equivalent to the action of the constancy mechanisms.  However,
a conclusion in favour of the second method would be too hasty.  Strict mathematical
analysis shows that a flawless conveyance of the geometry of perceptual space onto the
picture plane is impossible (see Appendix~2 for details).\footnote{%
  As a vivid example explaining the mathematical essence of the matter, we may cite
  the problem of creating geographical maps.  As is well known, a flat map cannot
  reproduce without distortion all that a globe conveys effortlessly: the shape and size
  of continents, and so on.  Therefore different cartographic projections exist, in each
  of which some elements~--- but not all~--- are reproduced without distortion, and
  there are no maps that correctly render everything.}
  More precisely, such an image can be correct only partially, and will contain, in
some elements, deviations from the geometric characteristics of the perceptual space
being depicted.

From this two conclusions may be drawn.  First, the creation of a ``geometrically ideal
image'' is in principle impossible.  The exceptional importance of this conclusion should
be emphasized: precisely this impossibility makes the search for approximations to the
ideal inevitable, and different methods of spatial construction~--- sometimes very
different from one another~--- naturally arise.  Second, it is evident that the depiction of
nearby and shallow spaces by the first method, where the depth cues that cannot be
reproduced on the canvas play the main role in the visual perception of the external
space, is entirely excluded.  In this case only the second method directly conveys the
aggregate effect of visual perception, accounting for all geometric transformations,
including those connected with the irreproducible depth cues.

\srcpage{55}

The broad comparison of the two possible systems of depicting perceptual space on a
canvas~--- while essentially correct~--- does not permit a definitive choice, and a more
detailed examination of their properties follows below.

We shall retain for both methods of depicting perceptual space the term \emph{systems
of perspective}.  Both systems will equally be scientific, since they are based on the
same regularities of visual perception and mathematics.  For the first of the named
systems it is necessary to know how to construct the retinal image; it coincides with
the central projection of spatial objects onto a plane.  Such projections are readily
obtained by the rules of linear perspective, and the corresponding perspective system
is called \emph{linear}.  (In mathematics the term ``linear'' is equivalent to
``rectilinear.'' It was applied in due course to emphasize that the image is obtained by
central projection along straight lines.)\footnote{%
  A more precise idea of this system is given by the expanded name: ``central
  rectilinear linear perspective.'' Wherever no misunderstanding arises, the standard
  abbreviated name ``linear perspective'' will be used.}
  This perspective type evidently coincides with that which has been taught in all art
schools for centuries and is usually called simply perspective or, as though to mark
its exclusivity, ``scientific perspective.''  The notion of the exclusivity of linear
perspective is, of course, wrong, since the second type of perspective is no less
effective and no less scientific.  The perspective system corresponding to the second
of the methods of depicting space considered above will be called \emph{perceptual}.
This is evidently connected with the fact that in it the geometry of perceptual space is
conveyed directly.

Below, the concept of \emph{axonometry} (parallel perspective) will frequently arise.
Axonometry will denote a system of perspective construction in which the property of
parallelism of straight lines is preserved.  Axonometry is a special case of both the
linear and the perceptual system of perspective.  It is well known that when depicting
small, simultaneously highly remote objects (theoretically: infinitely remote ones) in
the system of linear perspective, they may be depicted by the rules of axonometry.
Appearing in the system of linear perspective when depicting very distant planes,
axonometry for that very reason plays no special role in that system.  In the system of
perceptual perspective, analogous axonometric constructions apply for distant planes but
also, additionally, when depicting nearby space (as will be shown below).  For this
reason axonometry plays a very important role in the theory of the system of perceptual
perspective.  These questions are discussed in more detail in Appendix~5.\footnote{%
  Natural visual perception is exactly axonometric at two limiting distances, and,
  moreover, is more axonometric than linear perspective everywhere in between (see
  Appendix~4).}

\srcpage{56}

The system of linear perspective is well known; let us confine ourselves here to the
briefest statement of its geometric properties.  The basis of this system is the method
of central projection from some fixed point (in which the eye of the viewer is imagined
to be located) onto a plane perpendicular to the main line of sight.  Thus the basis of
linear perspective is monocularity and immobility of the viewpoint.  Since this same
construction scheme is characteristic of cameras and other analogous devices, linear
perspective may with full justification also be called optical or photographic perspective.

The laws governing image construction in such a perspective system are elementary.
Let us recall only that it has a number of simple geometric properties: straight lines of
objective space are depicted as straight lines; a group of horizontal parallel straight
lines lying in objective space is depicted as lines having a common vanishing point on
the horizon line; the sizes of depicted objects decrease with increasing depth in a simple
proportional relationship; and so on.  The well-known shortcomings of the system of
linear perspective have called forth a whole series of its ``improved'' modifications
permitting a larger angle of view, and so on.  However, all these modifications reduce
to projecting points of objective space in one manner or another onto some surface.  In
essence all these varieties are of one type, and for sufficiently small angles all coincide
with the ordinary system of linear perspective.  Therefore below only the latter will be
discussed, as representative of all systems of the same type.

As has been emphasized more than once, a person does not see the optical projection of
external space that has formed on his retina.  Therefore an image constructed by the
strict rules of linear perspective may in itself not correspond at all to visual perception.
As noted above, such correspondence would arise if, when contemplating the picture, the
viewer's visual perception system were subconsciously to carry out processes equivalent
to the action of the constancy mechanisms activated by contemplating objective external
space.  As a result of these subconscious processes, the retinal image would as it were
``stretch'' in the required regions and required directions.  The stretching can be of two
kinds~--- uniform and non-uniform.  The former, increasing sizes while not altering the
geometric form of figures, is characteristic of distant regions of space; the latter, in
which sizes change little but the form of figures is sharply transformed, is characteristic
of nearby regions of space.  (This question is discussed in more detail in Appendix~3.)

\srcpage{57}

The system of perceptual perspective is connected with the direct conveyance on the
canvas of the geometric properties of perceptual space.  Since perceptual space is built
by the full set of processes connected with visual perception, in the system of perspective
construction under discussion one can and should account for both the strong distortions
of the forms of the retinal image for nearby regions of space, and the strong ``uniform
stretching'' for its distant regions.

To systematise the study of the properties of this method of constructing a perspective
image, let us first consider a special case of it, in which the assumption of monocularity
is retained and the action of the form-constancy mechanism is not accounted for.  If,
based on experiments in the psychology of visual perception, the action of the
size-constancy mechanism is expressed in corresponding mathematical equations, it
becomes possible to depict not only the retinal image (to which linear perspective
corresponds) but also to construct the image that arises after the retinal image has been
subconsciously processed by the brain (perceptual perspective).  (The mathematical
aspects of this problem are presented in Appendix~3.)

\figblock{17}{Road and mountains in the systems of linear perspective (left) and
  perceptual perspective (right).  A horizontal surface covered, for clarity, with square
  tiles; mountains on the horizon.  Solid lines: the field of sharp central vision; dashed
  lines: peripheral vision.  In the perceptual system the tiles near the viewer are
  axonometric (parallel edges remain parallel), while distant regions are expanded
  uniformly relative to the linear-perspective image.}

Let us give the schematic image of a horizontal surface (for clarity imagined covered
with square tiles) in both systems (Fig.\,17).  The left diagram shows a horizontal field
with mountains on the horizon; the conventional image of the ground surface and the
mountains is divided into two regions~--- a central one shown by solid lines, and lateral
ones shown by dashes and without the ``chequered'' highlighting of individual tiles.

This division into two types of image is connected with the following circumstance.
The retina of the human eye is not homogeneous.  Its central (main) part gives high
acuity of perception; it corresponds to a comparatively small angle of vision whose size
depends on what is taken as ``sharp vision,'' while the large peripheral visual field
gives a very indistinct image, permitting ``sensing'' of the generalized outlines of
objects in objective space rather than examining them.  Provisionally, on Fig.\,17 the
region corresponding to the angle of sharp vision is given in solid lines, while the
peripheral-vision region is in dashes (the latter region is shown only partially).

\srcpage{58}

In the right diagram of Fig.\,17 the same subject is shown in the system of perceptual
perspective.  Comparison of the two images allows one to indicate a number of features
of perceptual perspective.  Let us first compare those regions corresponding to the angle
of sharp vision, i.e.\ conveyed by solid lines.  Comparing the far parts (adjacent to the
mountains, and the mountains themselves), one finds~--- as one would expect~--- that the
difference between the linear and perceptual systems reduces to a noticeable enlargement
in the latter, with preservation of geometric similarity.  One should at once note,
however, that this enlargement does not extend over the entire picture plane~--- for at
the base of the picture planes (it is assumed that the lower edge of the canvas coincides
with the base of the picture plane) the width of the square tiles in both pictures is the
same, i.e.\ the initial scale for nearby regions of space is the same in both images.

This circumstance explains the difficulties of artistic photography in the mountains.
When an alpinist photographs his companion against a background of mighty mountains,
after developing the film he sees him against a background of unremarkable modest
hills.  The whole point is that the alpinist's living visual perception ``stretched'' the
distant planes (the mountains), while the camera rendered the retinal image with
complete accuracy and without any transformation.  The unequal stretching of the near
and far planes leads, furthermore, to the fact that straight lines of objective space may
correspond to curved lines in the system of perceptual perspective.

The most interesting comparison is that of the depiction of near regions of space in the
two perspective systems.  In the system of linear perspective, nearby regions are depicted
by the same rules as distant ones, as is shown by the fact that both near and far sections
of the paths' edges share a common vanishing point on the horizon line.

The character of the depiction of nearby space in the system of perceptual perspective is
quite different: the image has become \emph{axonometric}~--- the edges of the square
tiles' images are shown by two pairs of parallel lines, i.e.\ parallels of objective space
are conveyed as parallels in the picture.  Thus axonometry (parallel perspective) proves
to be, as already noted, a special case of perceptual perspective for very nearby regions
of space.  Precisely because axonometry is connected with the natural visual perception
of the immediate surroundings, it plays a very important role in the work of artists who
convey nearby space.  The well-known axonometric character of distant planes in linear
perspective has limited artistic interest and is manifested, for example, when small
structures located on far planes are depicted.  Sometimes this axonometric limit
inherent in linear perspective is invoked to explain such depictions of the foreground,
claiming that the artist as it were mentally transports himself ``to infinity.''  But this
explanation would be contrived and formalistic, for no reasonable person, in order to
convey the appearance of a small object, would run to a considerable distance from it~---
he would rather come closer.

\srcpage{59}

It is therefore not surprising that for depicting large buildings and ensembles filling the
entire visual field, linear perspective is usually employed, while for architectural details
(i.e.\ objects viewed from close range) axonometry is invariably preferred.  This is not
only simpler in execution, but also more correct, i.e.\ closer to visual perception.

The natural appearance of axonometry is equally natural in all those cases where a
person directly conveys his visual perception of nearby space.  This is characteristic not
only of children's drawings, but of many artistic cultures of East and West.

In the right image of Fig.\,17 the region corresponding to the angle of sharp vision is
shown diverging in depth, which is connected with the ``stretching'' effect for those
parts of the retinal image corresponding to remote regions of space.  One should not
suppose that this is a result of some mathematical abstraction having no direct relation
to living visual perception.  If one lowers the gaze and looks carefully at the ground at
one's feet, mentally fixing the area of sharp vision, and then raises the head to look at
the horizon, one will have the sensation of an expansion of the angular measure of the
field of sharp vision~--- even though this angle remains in fact unchanged.

The transformations of the retinal image that resulted in the left diagram of Fig.\,17
acquiring the appearance shown in the right diagram are connected with the action of
the size-constancy mechanisms.  It is known, however, that only the field of sharp vision
is subjected to the transformations of these mechanisms.  The peripheral regions of the
retinal image undergo only the simplest transformation, explained by a ``scaling
alignment'' of the fields of sharp and unsharp vision and the preservation of image
continuity.  For this reason the field of peripheral vision is characterized by spatial
constructions close to the retinal image, i.e.\ close to linear perspective.  In the most
general terms this is shown by the dashed lines on the right diagram of Fig.\,17.

\srcpage{60}

In conclusion it should be mentioned that in fact a person is unable to take in with
sharp vision the entire region shown by solid lines in Fig.\,17.  In order to see space
from nearby regions to the horizon, one must shift one's gaze approximately three times,
gradually raising the head.  The provisional division into regions of sharp and unsharp
vision is given in the diagram only in the horizontal, not the vertical, direction.  In
practice a person will see sharply either one or two of the lower rows of black-and-white
tiles, or the horizon region and the tiles lying above the horizontal strip in which there
is no division into white and black tiles, or an intermediate region.  It may also be
noted that the curvature of lines running toward the vanishing point on the horizon will
not be conspicuous, since within each of those three regions, viewed separately, the
curvature in question is not very great.  For this reason the diagram in the figure is not
an example of a ``correct'' image, but merely a pretext for reasoning about the character
of human visual perception.

The striking difference in the visual impression of the images shown schematically in
the two perspective systems speaks, in particular, to the fact that the system of linear
perspective cannot fully satisfy the artist.  Outstanding masters of the visual arts,
sensing acutely the insufficient correspondence of images built in the system of linear
perspective to the visual perception of space, ``corrected'' it by inserting elements of
the system of perceptual perspective.  An analysis of the drawings of Bryullov, Polenov,
Vereshchagin, Repin, Serov, and other artists carried out by M.\,F.\,Fedorov showed
that the most typical deviations from the strict rules of linear perspective characteristic
of them can be reduced to three: smooth curving of lines that in reality are straight,
exaggeration of the sizes of objects in the far plane, and several different vanishing
points for objectively parallel lines.  All these deviations from strict linear perspective
are an obvious consequence of the artists' aspiration to approximate their images to the
visual perception of space, based on typical features of perceptual perspective.\footnote{%
  M.\,F.~Fedorov. \emph{Drawing and Perspective}. Moscow, 1960, p.\,63.  Such
  corrections to the system of linear perspective are of course not peculiar to Russian
  art alone.  Moreover, they are by no means limited to the goal of approximating visual
  perception, but often serve to intensify expressiveness.}

\srcpage{61}

Although the artists mentioned above introduced deviations from the strict system of
linear perspective~--- deviations understandable from the standpoint of the theory of
perceptual perspective~--- they firmly adhered to the latter as the basis for depicting
spaces extending in depth.

For distant portions of space, linear and perceptual perspective practically coincide;
this is visible in Fig.\,17 and follows from the properties of the process of transforming
the retinal image in the visual perception system noted above, in which the image of
distant regions of space undergoes uniform ``stretching'' with preservation of geometric
similarity.  To reveal the difference between the linear and perceptual systems of
perspective one must turn to the depiction of some simple object situated in the region
of the near foreground.  With this aim, let us examine various schemes for depicting an
interior~--- a shallow and very small rectangular room.

\figblock{18}{Possible variants for depicting a shallow, narrow interior.  Scheme~A:
  standard axonometric exterior view.  Scheme~B: partial axonometric interior (only one
  corner).  Scheme~C: axonometric view of the far wall only.  Scheme~D: two separate
  axonometric spaces with a gap (inadmissible discontinuity).  Schemes~E and~F:
  partial interiors omitting some walls.  Scheme~G: linear-perspective interior~---
  parallel lines are forced to converge, and the two bold lines become mutually
  perpendicular.}

Figure~18 gives several variants of such schemes.  Bold lines in all schemes show
lines drawn along the middle of the floor and of the left wall.  Both lines are parallel
to the corresponding edges of the parallelepiped forming the room.

Scheme~A is an ordinary axonometric image, but it depicts the room from outside,
whereas it is necessary to show the room from the standpoint of a person standing
inside.  If the same technique is used for this purpose, scheme~B arises, whose main
deficiency is a partial depiction of the interior; it does not permit, for example, both
walls of the room to be shown simultaneously, even though a person standing inside it
is capable of seeing these walls at once.  The most rigorous type of axonometric image
of the room when viewed from the front can be considered scheme~C, which shows the
far wall.  Since in axonometry the distance from one side wall to the other, or from
ceiling to floor, does not change with depth, the depicted width or height of the room
will be the same for all distances from the viewer.  This leads to the side walls, floor,
and ceiling being projected into lines, and the bold lines~--- into points (not shown on
the scheme).  Scheme~D is clearly inadmissible, since it contains discontinuities in the
image, but it correctly conveys two properties of the visible image of the room: the right
and left borders of the floor and the bold line drawn on it are not only parallel to one
another but also perpendicular to the far wall; the images of the left wall, right wall,
and ceiling have the same property.  Since a discontinuity in the image is inadmissible,
one may, just as in passing from scheme~A to scheme~B, proceed by partial depiction of
the interior~--- giving rise to schemes~E and~F.  The choice between the latter two
schemes is determined by the general artistic conception, since it is connected with the
``cutting off'' of planes secondary to the given composition.  Scheme~G gives an image
of the same interior in linear perspective (for definiteness it is assumed that the artist
stands at equal distances from the side walls, and the bold line on the left wall is drawn
at eye level).  The shortcoming of the last image is the violation of the clearly visible
parallelism of lines: the bold lines, which are in reality and in visual perception parallel
to one another, the artist is forced to depict as perpendicular lines.

\srcpage{62}

The schemes given here illustrate the assertion made earlier that the creation of an
``ideal'' picture is impossible.  Indeed, an ``ideal'' picture should have the following
properties: no discontinuities in the image; simultaneous visibility of the side walls,
ceiling, and floor; mutual parallelism of the edges formed by the intersections of the
side walls, ceiling, and floor (or at least only a slight deviation from this parallelism);
perpendicularity of these edges to the corresponding boundaries of the far wall (or at
most only a slight deviation).  All schemes shown are far from ideal; schemes~B, E,
and~F give only a partial interior; in scheme~B the aforementioned perpendicularity
property is also lost; scheme~D contains gaps; scheme~C shows no side walls, ceiling,
or floor; in scheme~G the parallel lines have lost this important property and the two
bold lines have even become mutually perpendicular.  As we see, both theoretically
possible systems of perspective can convey only some properties of the perceptual
spatial image, which brings up the problem of compensating for the distortions of the
visible image in them.  In accordance with the two-stage scheme of human visual
perception and the two perspective systems corresponding to those stages, two methods
arise of compensating for distortions~--- inevitable in any perspective system~--- of the
image created by the human visual perception system.

If the artist uses the system of linear perspective, then, as already noted in the previous
chapter, it is important that perception rely on all the depth cues.

Of the depth cues listed in the previous chapter, the artist, as already noted, can depict
only some of the monocular ones.  This immediately makes the attempt to compensate
for the distortions inherent in the images of nearby objects in the system of linear
perspective practically hopeless, since in contemplating nearby objects in reality it is
precisely the binocular depth cues that are the principal ones.  For this reason the artist
should not depict very close objects at all.\footnote{%
  When depicting such objects becomes necessary, artists deviate from the requirements
  of linear perspective.}
  If one turns to somewhat more distant regions of space, the role of the monocular
depth cues increases, and for distant regions they are the only ones.  In order to
increase the effectiveness of transforming the retinal image, the artist must emphasize
those depth cues that can be used for this purpose.  Such a cue as the diminution of
object sizes with depth is a direct result of applying linear perspective and therefore
requires no special attention.  Equally natural is the blocking of distant objects by
nearby ones, and the proximity of objects' images to the horizon line with increasing
depth.  These two cues are also a consequence of applying linear perspective.  A quite
different~--- non-geometric and therefore especially effective~--- character belongs to
the two other cues: aerial perspective, and systems of shadows on large surfaces.  A
not unimportant role in the ``deception'' of the visual perception system that must make
it feel the spatial authenticity of what is depicted in the picture is played by
``recognition,'' which in turn requires colours as close to nature as possible.

\srcpage{63}

However, in a picture painted in the system of linear perspective with all possible
(including non-geometric) depth cues employed, several depth cues cannot be
reproduced.  Moreover, a person contemplating the picture cannot fully abstract from
the texture of its surface, and so on.  Nevertheless, it is not the geometry of the image
itself (linear perspective) but precisely this work of the brain~--- analogous to the work
it performs in forming subjective three-dimensional space from the two-dimensional
retinal image~--- that gives the sense of depth characteristic of images built by the rules
of linear perspective.\footnote{%
  It is the reproduction of the depth cues (not this or that image geometry in itself)
  that provides the decisive contribution to the sense of depth of depicted space.  For
  this reason pictures conveying nearby regions of space are always more ``flat'' than,
  say, a landscape, since for nearby regions of space it is the binocular depth cues~---
  unreproduble in a picture~--- that are decisive.  How much we lose by being unable to
  reproduce these cues is vividly shown by stereoscopic and holographic images of
  nearby regions of space (while distant regions, under monocular observation and
  through a stereoscope, are practically identical).}
  In those cases where the artist does not follow the rules of linear perspective but
applies all five depth cues listed, the sense of spatial depth nonetheless arises.  This
is seen from the example of medieval Chinese landscape painting.

\figblock{19}{Xia Gui. \emph{Autumn Landscape.} Album leaf. Late 12th -- early 13th
  century.  A Chinese landscape employing depth cues (aerial perspective analogues,
  occlusion) without linear perspective, yet conveying a powerful sense of spatial
  depth.}

\srcpage{64}

The visual perception discussed here possesses a certain property of integrality.  The
brain of the viewer transforms the entire picture as a whole, finding some overall
approach to the spatial interpretation of the depicted subject.  This permits the required
transformation of the images of relatively nearby objects to be reinforced, using the
depth cues effective for distant objects (for instance aerial perspective), if those distant
objects are depicted on the same picture.

From childhood onwards the modern person deals mainly with images using the system
of linear perspective.  Photographs and films do so with complete strictness, and
television images do so in their main outlines.  Furthermore, for several centuries now
pictures, book illustrations, and the like have been based on the system of linear
perspective.  However much it distorts the visible images, our brain has grown
accustomed to these distortions and knows how to interpret~--- and hence ``not
notice''~--- them.  Facilitating this is also the rigidly programmed character of the
distortions, conditioned by the unambiguous rules of constructing linear perspective.
Moreover, the completely unnatural distortions of visible images bestowed on us by
photography (for example, cannon barrels photographed from close range from the
muzzle end, showing a hypertrophically enlarged muzzle diameter compared with the
more distant parts of the barrel) have gradually been transferred to artists' canvases
without provoking feelings of protest in viewers.  Yet the photographic image is far
from ``natural.''  This is indicated by observations of psychologists who noted that
peoples at a lower stage of civilization do not perceive photographs as a natural image
of the external world.

\srcpage{65}

It is difficult for us to imagine the enormous work of ``educating'' the viewer that the
masters of the Renaissance carried out.  Traces of this work can still be seen on their
canvases.  Renaissance paintings are often characterized by an emphasized depiction of
depth: long corridors, open windows or doors in which roads running into the distance
are masterfully depicted, in which chains of mountains with castles are seen~---
spatially separated by the skilful use of aerial perspective~--- and so on.  In these
canvases the artists were primarily solving the artistic problems before them, but it is
very likely that simultaneously they (intuitively, of course) strove to evoke in the viewer
emotions having an effect equivalent to the size-constancy mechanism.  It was noted
somewhat above that the viewer perceives and processes the picture integrally in his
consciousness as something whole; consequently, the sense of depth, evoked by the
emphasized depiction of the continuous receding vista, stimulated the processing also of
closer planes by the same laws, and intensified the artistic effect of applying linear
perspective for those regions of space where it already noticeably diverges from the
perceptual image.

It is also appropriate to recall that distant regions of space are depicted in the linear and
perceptual systems of perspective in geometrically similar ways.  Therefore linear
perspective, applied to distant regions of space, could not seem strange to a person
accustomed to perceptual perspective.  If one further takes into account that for distant
regions of space only the monocular depth cues~--- fully reproducible in a picture~---
are operative, then the stimulation of those centres of the human visual perception
system responsible for transforming the retinal image into three-dimensional perceptual
space (i.e.\ the stimulation of the sense of depth) should not have presented great
difficulties.  Gradually, habit with linear perspective made unnecessary such indirect
techniques for evoking the sense of depth in the depiction of nearby planes.

Depiction in perceptual perspective (this concerns, of course, nearby regions of space,
where the two perspective systems differ) is constructed taking into account the action
of the constancy mechanisms.  Here a repeated reworking of what is depicted by the
human visual perception system can lead to distortions of the visible image, and
therefore depth cues must be introduced into perceptual perspective with caution.

\srcpage{66}

By reason of the fundamental impossibility of conveying the image of perceptual space
on a picture plane without distortions, the artist is compelled to strive to completely
eliminate distortions in the depiction of the principal elements and to transfer them to
the secondary ones.

Let us examine in more detail the paths of compensation of distortions in the system of
perceptual perspective outlined above.  If only distant regions of space are depicted on
the canvas (a landscape whose foreground is sufficiently remote), then all depth cues are
fully appropriate here.  Since (as already noted) for distant regions of space linear and
perceptual perspective give geometrically similar images~--- images containing no
geometric distortions~--- there need be no difference in the methods of stimulating the
sense of depth (in particular by using aerial perspective).  This is seen from the
already-cited example of medieval Chinese landscape painting, in which the technique
equivalent to aerial perspective was widely used long before the Renaissance.

If the picture depicts distant regions of space together with nearby ones, the problem
becomes more complex.  To illustrate the resulting difficulties, let us turn to an example
(Fig.\,20).  Its upper part gives a schematic image of a road leading to distant mountains,
with two pedestrian paths at the sides, depicted in linear perspective.  The middle part
of the figure shows the same road in an image close to perceptual perspective~---
the far plane is given in linear perspective, and the near plane in axonometry (the kinks
in the lines arising from this simplified combination of the two image types testify to the
already-mentioned difficulty of coordinating axonometry and linear perspective on one
canvas).  For such an image to be correctly perceived, the action of the constancy
mechanism must be excluded, since it has already been approximately accounted for
(the image of the road is of the same type as the right image of Fig.\,17: linear
perspective passes into axonometry in the foreground).  In reality, however, the
appearance of the receding part of the road gives a sense of depth and provokes the
viewer, who begins to subconsciously transform the entire arising retinal image by laws
equivalent to the action of the size-constancy mechanism.  As a result, the near portion
of the road seems inclined~--- instead of a level horizontal road (the upper part of the
figure) an image of a road that first rises up a slope and only then becomes horizontal
arises.  This transformation of a level road into one that first goes uphill is the
distortion of the true picture, connected with the fact that when perceptual perspective
was used, the transformation of the picture as a whole by the human perception system
according to laws equivalent to the action of the size-constancy mechanism~--- activated
by the image of the far part of the road~--- was not excluded.

\srcpage{67}

Exactly the same phenomenon can be seen in the right image of Fig.\,17, where the
square tiles closest to the viewer tend to ``stand up vertically.''  This effect is observed
not only in theoretical schemes, but also in the canvases of artists.  The foreground of
V.\,Bednov's painting \emph{Bleaching of Linens} (Fig.\,21) is close to an image by
the laws of perceptual perspective: the width of the linens increases insufficiently
energetically for linear perspective as one moves toward the lower part of the canvas.
As a result, the meadow depicted in the foreground seems to lie on a hillside.  This
should not be understood as a criticism of the artist~--- such a depiction was probably
part of his artistic design~--- but it shows that the effect described above is not a
theoretical abstraction.

\figblock{20}{A road leading to distant mountains (with pedestrian paths) depicted in
  three ways.  Upper: linear perspective.  Middle: a hybrid~--- far plane in linear
  perspective, near plane in axonometric~--- showing the junction kink and the
  perceptual illusion that the near road rises uphill.  Lower: the discrete-axonometry
  solution~--- space divided into several independent axonometric zones, each masked at
  its junction (e.g.\ by a wall), avoiding both the junction kink and the unwanted
  depth-cue trigger.}

\figblock{21}{V.\,Bednov. \emph{Bleaching of Linens.}  The foreground is rendered
  close to perceptual perspective (width of linens increases too slowly for linear
  perspective), causing the meadow to appear to slope uphill~--- an example of
  unwanted re-transformation by the constancy mechanism.}

With the aim of suppressing the unwanted repeated transformation of an image made by
the rules of perceptual perspective by the human visual perception system, various
approaches are possible.  The simplest must be recognized as the method in which a
close foreground and a far plane are never depicted on one canvas.  In that case the
shallow nearby foreground, depicted by the laws of axonometry, though conveying the
volume of objects, does not stimulate effects equivalent to the action of the size-constancy
mechanism~--- precisely by reason of the small depth of the depicted space.  With such
separate depiction of an axonometric foreground and a distant plane, the problem of
coordinating them on one canvas~--- discussed above~--- also does not arise.

In those cases where it is necessary to show simultaneously the near plane and the
distance in perceptual perspective, and at the same time not to provoke the viewer into
a repeated transformation of the contemplated canvas by the size-constancy mechanism,
one may propose the method shown in the lower part of Fig.\,20.  The entire space is
divided into a series of (in this case two) discrete planes, each of which is depicted by
the rules of axonometry, and the junction sections of two axonometries are ``masked'' by
the image of a wall.  Here obvious use is made of the circumstance that in the system
of perceptual perspective, images of both shallow nearby and shallow distant regions of
space are axonometric.  Thus the described method reduces to decreasing~---
stepwise, discretely~--- the sizes of objects that diminish gradually toward the horizon
line (here: the width of the road), and, if necessary, ``masking'' the junctions that arise.
In such a picture construction (fitting entirely within the system of perceptual
perspective) the viewer perceives it as a series of successive axonometric spaces, and
the step-change in sizes does not merge into a sense of the continuity of the distance
that impels the person looking at the picture to its subconscious transformation according
to laws approaching the size-constancy mechanism.  At the same time the question of
coordinating linear perspective and axonometry on one canvas is also resolved.

In those cases where the step change in the scales of neighbouring axonometric spaces
is masked (for example by a wall) or justified by some other real cause, one can enhance
the sense of depth of the depicted space, since a certain justified gap arises between
neighbouring axonometries.  If this is not done, the distant axonometry will be perceived
as a kind of flat theatrical backdrop immediately adjoining the main axonometric space.

\srcpage{68}

What methods exist for compensating for the inevitably geometric distortions of a
sufficiently complete perspective image (violations of the apparent parallelism of lines,
of the relative sizes of objects, etc.)?  The circle of questions connected with this will
be illustrated by the already-discussed example of the image of a small and shallow
room.

Let us give two images of the room (Fig.\,22), the upper of which, built by the rules of
linear perspective, may, with the qualifications already given, be considered the retinal
image.  If one compares it with the visible image, one will easily note a number of
discrepancies: as a result of the action of the size-constancy mechanism, the visible
width of the far wall~$AB$ will be very close to~$CD$, angle~$\beta$ will become
almost a right angle, and the floor-boards will practically acquire the configuration of
rectangles.  Let us attempt to render this on a flat sheet of paper.  The result will be
the image given in the lower part of Fig.\,22.  It is clear that a more correct rendering
of the perceptual image of the far wall, floor, and ceiling occurred at the expense of an
even stronger~--- than in linear perspective~--- distortion of the visible image of the side
walls, which were turned into barely perceptible trapezoids.  It is easy to see that all
the reasoning given about the floor and the width of the far wall applies equally to the
side walls and the height of the far wall: the size-constancy mechanism should nearly
equalize the segments~$BE$ and~$DF$, and angle~$\alpha$ should become close to a
right angle.  Yet in the lower image the angle designated~$\alpha$ became even more
obtuse, and the relative heights corresponding to~$BE$ and~$DF$ converged
insufficiently.  Thus, not only is the impossibility of adequately depicting on a plane the
spatial perceptual image illustrated once more, but~--- more importantly~--- the method
of compensating for perspective distortions in the system of perceptual perspective is
shown.

\figblock{22}{Method of local compensation of distortions in the system of
  perceptual perspective.  Upper: linear-perspective image of a small room (the retinal
  image).  Lower: perceptual-perspective image of the same room~--- the far wall and
  floor are rendered correctly at the cost of the side walls becoming barely perceptible
  trapezoids (the ``law of conservation of distortions'').}

\srcpage{69}

In the system of perceptual perspective, the correct rendering of some elements of the
image (for instance the floor) occurs at the cost of distorting other elements of the same
image (for instance the side walls); if these distortions seem to the artist too strong,
he is entitled (when possible) simply not to depict the corresponding elements.  In
particular, in the lower drawing of Fig.\,22 one could have not shown the side walls at
all and used a scheme of the type~18E.\footnote{%
  Depicting a room without side walls should not be regarded as a gross deviation
  from artistic truth.  When a person concentrates attention on the floor and far wall
  surfaces, the side walls fall in the field of indistinct peripheral vision; their presence
  is sensed rather than registered in detail.  Therefore the absence of a depiction of
  the side walls can be more truthful than their detailed and, moreover, strongly
  distorted rendering in a picture.}

\srcpage{70}

In principle another approach is also conceivable, in which the question of the general
properties of the image is settled before the artist sits down to work.  For instance, it
is decided that errors are permissible for verticals but excluded for horizontals.  Then
the corresponding ``compensations'' and the distortions inevitably connected with them
would be built into the arrangement of points of the ``empty'' space as it were in
advance, and, just as the system of linear perspective, the system of perceptual
perspective would become ``rigid''~--- with rules of construction formulated in advance.
Then any point of objective space would map to one and the same point on the picture
plane regardless of what object is at that point of the objective space, which would
completely eliminate the artist's ``arbitrariness.''  Of course, such a rigid system of
perceptual perspective~--- mathematically feasible~--- is never used by artists, but it is
convenient to introduce when carrying out art-historical analysis.  (An example of using
such an approach will be given in Chapter~X\@.)  Thus theoretically the following
variants of perspective systems are possible:

\begin{enumerate}
\item The system of linear perspective (always rigid);
\item The system of perceptual perspective
  \begin{enumerate}
  \item free,
  \item rigid.
  \end{enumerate}
\end{enumerate}

To conclude the general survey of perspective systems, two special questions must be
addressed.  The first is the role of peripheral vision, and the second is that of the
form-constancy mechanism.  It was noted above that the central field of sharp vision
and the wide field of peripheral vision have different properties, and this circumstance
must somehow affect the work of artists.

Suppose a shallow interior must be depicted, but its width exceeds the field of sharp
vision.  Then in the consciousness of a person standing equidistant from the side walls
and looking straight ahead, the picture schematically represented in Fig.\,23A will
arise.  Assuming for simplicity of reasoning full constancy of perception at the small
distance to the far wall (this condition is approximately met in practice), one may assert
that the person will see the floor-boards in the central part of the room as a group of
parallel lines simultaneously perpendicular to the far boundary of the floor.  As for the
side walls of the room, they will fall in the field of indistinct peripheral vision (this
is shown in the figure by dashed lines), and since the constancy mechanisms do not
transform the peripheral vision regions according to their laws, the room's walls will be
seen in indistinct direct perspective, close to ordinary linear perspective.  As for the
floor-boards near the side walls, their image will also be ``blurred,'' and since they are
much smaller than the walls, their configuration will be substantially less definite than
that of the walls~--- a fact marked in the scheme by the indistinctly seen floor-boards
simply not being depicted.

\srcpage{71}

Wishing to clarify the appearance of the side walls and the adjacent space, the artist
will shift his gaze, say, to the left.  Then the field of sharp vision will be at the left
wall, and in his consciousness the picture shown in scheme~23B will arise.  He will now
clearly see the left floor-boards also as a series of parallel lines~--- since the field of
sharp vision is transformed by the constancy mechanism~--- simultaneously parallel to the
left boundary of the floor.  A completely analogous situation will arise when the gaze is
shifted to the right.  If the artist now wishes to convey his vision in a picture, depicting
both the central and lateral parts of the floor as he sees them, this will lead to scheme~23C.

\figblock{23}{Depiction of the floor of a wide room, requiring corrective shifts of gaze.
  Scheme~A: looking straight ahead; the central floor-boards appear as parallel lines
  perpendicular to the far wall; side walls are in indistinct linear perspective.
  Scheme~B: after shifting gaze to the left wall; the left floor-boards also appear as
  parallel lines, but now parallel to the left wall.  Scheme~C: combining the views
  produces three mutually incompatible sets of parallel lines, with ``empty triangles''
  at their junctions.}

The composite scheme~C vividly illustrates the difficulty the artist will face.  If one
simply follows it, it becomes entirely unclear how to ``join'' the images of the side and
central groups of floor-boards: these three series of parallel lines in the region of
``joining'' will come into glaring contradiction with one another, forming ``empty''
triangles in the picture.  Anticipating what is to come, we note that this problem
continually troubled artists.  In some epochs it was preferred not to show the side walls
at all; sometimes they were shown while omitting or ``masking'' the centre; sometimes
only one corner of the room was shown, as in scheme~18B.  In the system of linear
perspective the image distortions were distributed ``equally,'' giving everywhere the
same floor-boards converging toward the far wall~--- and thereby everywhere violating
visual perception, since wherever in reality a person looks, he sees the floor-boards with
parallel edges in that direction.

\srcpage{72}

The solution of the problem of depicting the floor in full accord with visual perception
proved to lie within the reach of Post-Impressionism.  Let us turn in this connection to
the well-known canvas of Van~Gogh, \emph{The Bedroom in Arles} (Fig.\,24).  Van~Gogh
solved this problem with truly remarkable simplicity.  He conveyed the floor structure
by a series of seemingly disordered quadrilateral figures~--- as if images of rectangular
elements of a parquet floor.  Yet this ``disorder'' has a hidden rational structure.
Human perception has the property of ``uniting'' even amorphous spots into lines and
figures; therefore when looking at Van~Gogh's canvas one clearly senses that in the
central part of the room the floor structure is formed by nearly parallel strips running
toward the wall with the window.  Van~Gogh made them not strictly parallel simply
because he saw them that way (strict parallelism in Fig.\,23 was introduced for
simplicity of reasoning).  But remarkably, even when shifting the gaze toward the side
walls one again sees an approximately parallel structure~--- just as in scheme~23C,
but without the troublesome empty triangles, since the ordering is not exact.

\figblock{24}{Vincent van~Gogh. \emph{The Bedroom in Arles.}  Van~Gogh's
  perceptual solution: the floor is rendered as seemingly disordered quadrilaterals that
  the viewer integrates, from any gaze direction, into approximately parallel strips.
  This avoids both the ``empty triangles'' of scheme~23C and the forced convergence
  of linear perspective.}

Let us analyse with this goal the problem of depicting a spatial corner~--- the junction
of ceiling and two walls~--- in the system of perceptual perspective.  Since a corner is
not a linear magnitude, it is transformed not by the size-constancy mechanism but by
the form-constancy mechanism.

Consider the spatial corner at one vertex of a parallelepiped, and let us make no
distinction between looking at this corner ``from inside'' or ``from outside,'' since
the line of reasoning does not change.  Figure~25 shows a certain vertex~$A$ of the
parallelepiped; by the well-known properties of a parallelepiped,
$\alpha = \beta = \gamma = 90°$, from which $\alpha + \beta + \gamma = 270°$.  The
image of these same three angles on a plane, completely independently of the chosen
perspective system, will always give $\alpha + \beta + \gamma = 360°$, since on the
picture plane these three angles fill the entire surroundings of point~$A$.  The retinal
image has exactly the same character: in it too the sum of these three angles will
be~$360°$.  As is well known, the form-constancy mechanism brings the retinal image
closer to the true form of an object if it is well known to the observer.  Since we all
live in a ``right-angled'' world~--- that is, most often deal with right angles (recall the
configuration of houses, rooms, wardrobes, tables, boxes, books, etc.)~--- the true
magnitudes of angles~$\alpha$, $\beta$, and~$\gamma$ are absolutely clear to the
observer at first glance at the parallelepiped, even if he has not studied geometry and
has no idea of the geometric meaning of the concept ``right angle.''

\srcpage{73}

The form-constancy mechanism will therefore transform the retinal image in the
direction of bringing it closer to the true values of the three angles in question; however,
it is not capable of carrying out so deep a transformation that the angles would be seen
as right angles.  Indeed, when looking at a round table standing in a room, a person,
though seeing it as less oval than a photograph would show it, can never see it as an
exact circle.  Therefore the human perception system will yield
$270° < \alpha + \beta + \gamma < 360°$~--- which it is fundamentally impossible to
depict on a plane, regardless of which perspective system is applied, since on a plane
this sum is always exactly~$360°$.  Experiments carried out by the author showed that
most observers estimate the sum of these three angles at around
$300°\text{--}320°$.\footnote{%
  In such experiments observers should be asked to estimate the apparent magnitude
  of each of the three angles separately, as if trying to draw it on three separate sheets
  of paper, and only then to sum these magnitudes.}

\figblock{25}{Spatial corner at vertex~$A$ of a parallelepiped.  Objectively
  $\alpha = \beta = \gamma = 90°$ and $\alpha + \beta + \gamma = 270°$.  On any
  planar image, whether retinal or on canvas, the sum equals~$360°$.  The
  form-constancy mechanism brings the perceived sum toward~$270°$, but cannot
  reach it; experiments give approximately $300°\text{--}320°$.}

The crux of the difficulties connected with direct depiction of the visible image in the
system of perceptual perspective reduces to the fact that the visible image perceived by
a person is three-dimensional.  The sum $\alpha + \beta + \gamma < 360°$ is possible
only for spatial objects, and precisely this inequality creates in a person the sense of
spatiality~--- and it is fatally just this that cannot be depicted.  True, the actually seen
angles, for which this sum is of the order of~$300°\text{--}320°$ as noted, correspond
to some spatial trihedral angle considerably ``flatter'' than the angle of a parallelepiped.
The visible image is both spatial and at the same time approximated to a flat one.  It is
probably the latter that simplifies the depiction (with the inevitable distortions in any
perspective system) of space on a plane.

The spatiality of the perceptual image is the main cause of the difficulties connected
with depicting these images on a plane.  It is well known~--- and this is not the place
to go into the corresponding mathematical details~--- that an arbitrary spatial surface
(which one may visualize as a corresponding spatial paper model) cannot be
``flattened,'' laid for instance on a flat table, without the formation of tears or folds.
Therefore attempts to give an absolutely complete perceptual image inevitably lead to
individual portions of the depicted surfaces ``riding over'' one another, or to gaps
forming in the image.  Since neither is permissible, the artist prefers local distortions
of the visible picture.

\srcpage{74}

Another question that arises when accounting for the action of the form-constancy
mechanism is the dependence of the properties of perceptual space on the objects
disposed in it.  Let us illustrate this with a simple example.

\figblock{26}{Illustration of the deformation of perceptual space by the form-constancy
  mechanism.  Upper: four points~$A$, $B$, $C$, $D$ (coloured balls placed at the
  corners of a horizontal square) as they appear in perceptual space when their true
  arrangement is unknown~--- as a near-parallelogram.  Lower: when the same four
  points are perceived as the corners of a familiar square stool seat, the form-constancy
  mechanism deforms the near-parallelogram toward a square; $C$ and $D$ shift to new
  positions~$C'$ and~$D'$.}

Suppose that in real space there are four points, marked by coloured balls, and suppose
that the person viewing them from a short distance does not know their actual
arrangement in objective space relative to one another.  Then in perceptual space their
mutual arrangement will correspond to the retinal image transformed only by the
size-constancy mechanism, since the brain will have no supplementary information about
them allowing the observed configuration to be brought closer to the true one in
accordance with the form-constancy mechanism.  Figure~26 shows four points~$A$,
$B$, $C$, $D$, arranged at the corners of an imaginary horizontal square, in the form
in which they are reflected in perceptual space.  If now, without moving these points
(balls) in real space, one arranges for them to correspond to the corners of the square
seat of a stool familiar to the viewer, then~--- although in the retinal image they will not
shift~--- in perceptual space the near-parallelogram figure~$ABCD$ will be deformed
in the direction of being brought closer to a square, for example acquiring the form
$A_1B_1C_1D_1$.  Let us suppose for simplicity of reasoning that points~$A$ and~$A_1$
and correspondingly~$B$ and~$B_1$ coincide in perceptual space.  Then one may speak
of points~$C$ and~$D$ having shifted to positions~$C_1$ and~$D_1$ under the action
of the form-constancy mechanism, which in turn was able to manifest its action only
because the observer saw the coloured balls lying at the corners of a well-known object
having a square upper face~--- an object about which the observer had \emph{a priori}
information.

\srcpage{75}

This example shows that perceptual space lacks the most important property of
objective space~--- that each point of space is immobile relative to other points of the
same space regardless of what objects are at those points at a given moment.  Perceptual
space is characterized by the variability of the mutual positions of points depending on
the objects situated at them, or more precisely, depending on the \emph{a priori}
information about those objects available to the observer.

The phenomenon described here~--- by which objects can ``deform'' perceptual space~---
creates a fundamental difficulty in constructing a theory of perceptual perspective,
since any one-to-one mapping of ``empty'' perceptual space onto some plane will
immediately change as soon as that space begins to be filled with objects familiar to the
viewer.  It is only the independence of the properties of real space from the objects in
it that permits the creation of a unique and simple system of linear perspective
(unique in the sense that each point of space corresponds to one definite point on the
picture plane~--- regardless of what is placed at that point of real space).  An additional
difficulty in the attempt to create a rigorous system of perceptual perspective is that its
very ambiguity (the degree of deformation of perceptual space by images of objects)
depends on the \emph{a priori} information about the objects available to the viewer.
Naturally, the character of the \emph{a priori} information available to one person may
differ from that available to another, and therefore the degree of deformation of the
``empty'' perceptual space will be different in these two cases.

In order not to be constrained by the difficulties described above, in Appendix~2
\emph{perceptual perspective in the narrow sense} is defined as that which takes into
account the action of the size-constancy mechanism and does not take into account the
action of the form-constancy mechanism.  The reasonableness of this approach is now
clear: the size-constancy mechanism acts on the images of objects depending on their
distance from the viewer, not depending on \emph{a priori} information about them.
Therefore the size-constancy mechanism can operate even when observing a space
``filled'' with mathematical points, and its effectiveness will not change when real
objects appear in that ``empty'' space.  In this approach, accounting for the action of the
form-constancy mechanism must be carried out separately, in relation to the local
regions of space associated with the objects situated in them.  This is precisely the
method adopted in the present book.

\srcpage{76}

Thus the system of perceptual perspective has proved to be quite flexible.  After the
retinal image has been transformed by the visual perception system and the artist
proceeds to transfer this perception to the picture plane, he is confronted with a
fundamental property of this process~--- the impossibility in principle of an exactly
accurate transfer of the geometric properties of perception to the picture plane.  This
impossibility is connected with the action of both the size-constancy mechanism and the
form-constancy mechanism.  Depending on how the artist resolves the problem of
compensating for the inevitable distortions~--- what he decides to convey exactly and
what to ``distort''~--- different types of perceptual perspective may arise, and this
diversity of types is indeed observed in the history of art.  In the present book two
types of such systems will be examined~--- those connected with the art of the Middle
Ages and with the art of Post-Impressionism~--- but their number is of course larger.

In conclusion, a comparison of the properties of the two principal systems of
perspective~--- linear and perceptual~--- must be given.

\medskip
\noindent\textbf{1.}~When distant regions of space are depicted in isolation, both
systems coincide and both accurately reproduce natural visual perception.  Their
difference appears when it becomes necessary to depict the middle, and still more the
near, foreground.

\medskip
\noindent\textbf{2.}~Regardless of the perspective system used, depiction of the middle
and near foreground without deviations from the geometry of visual perception (without
``errors'') is impossible.

\medskip
\noindent\textbf{3.}~The depiction of nearby regions of space in the system of linear
perspective is ineffective, since the binocular depth cues~--- which play the principal
role in contemplating these regions and which cannot be reproduced in a picture~---
dominate.  The resulting image will be perceived by the viewer as containing an
especially large number of very gross ``errors.''  Perceptual perspective, in which the
geometry of visual perception is depicted, knows no such limitations (hence artists
painting a group portrait with equally-sized heads~--- without their perspective
diminution~--- are in practice working within the system of perceptual perspective).

\srcpage{77}

\medskip
\noindent\textbf{4.}~The rules of image construction in the system of linear perspective
are determined by the unambiguous laws of optics, i.e.\ are not connected with the
working of human consciousness (a photograph is the paradigm).  The construction of
an image in the system of perceptual perspective~--- i.e.\ the conveyance on the picture
plane of the geometry of visual perception~--- cannot be accomplished without relying
on the working of human consciousness.  This reliance reduces not only to the
subconscious work of the perception system, but also to the conscious choice of those
elements of the image that may permissibly be distorted, and those that should be
conveyed in exact accordance with visual perception.  Since this choice is made by the
artist in accordance with the artistic tasks he is solving, many geometrically different
depictions of one and the same pictorial space, painted from one and the same viewpoint,
can exist.  Consequently, the system of linear perspective is unambiguous in character,
while the system of perceptual perspective is multivalent.

\medskip
\noindent\textbf{5.}~The inevitability of distortions in the geometry of visual perception
when conveying the middle and near foreground brings with it the necessity of
compensating for these distortions.  In the system of linear perspective, where the
distortions are rigidly fixed by the rules of its construction, the only method of
compensation is the emphasized depiction of those depth cues that the artist can
reproduce (in order to stimulate those centres of the visual perception system that
transform the retinal image into perceptual space).  Since the artist can reproduce only
some of these cues in the picture, the compensation will be incomplete.  In the system
of perceptual perspective, the artist, conveying exactly the geometry of those elements
that are important to him, permits a distorted image of others.  A ``law of conservation
of distortions'' operates here: the correct image of one element leads to the intensified
distortion of another.  Sometimes this threatens distortions so severe that the artist
prefers simply not to depict the corresponding element falling in the field of vision (for
example, omitting the side walls of an interior).  In the system of linear perspective,
as a consequence of the inadequacy of the image built by its rules to natural visual
perception, the distortions are distributed uniformly among all elements (i.e.\ everything
is depicted ``incorrectly''), while in the system of perceptual perspective the ensemble
of locally correct elements will be supplemented by a corresponding number of elements
distorted more severely than in the system of linear perspective.

\srcpage{78}

The properties of the two perspective systems listed here give no grounds for regarding
one as straightforwardly superior to the other.  Rather, they should have different
domains of application.

The comparison of the two perspective systems was begun with the observation that for
the isolated depiction of distant regions of space they coincide.  This permits the
formulation of an alternative approach to the problem.  One may assert that only the
perceptual system of perspective exists, and the linear is its special case for distant
regions of space.  This formulation is based, in particular, on the fact that~--- unlike the
system of linear perspective~--- the perceptual system can give images from the very
nearest to arbitrarily distant regions of space.  A more detailed exposition of this
formulation will rationally be undertaken after the features of the system of perceptual
perspective have been studied and clarified in sufficient detail through examples from
works of art.

The system of linear perspective has blinded many with the mathematical formulas and
optical-geometric constructions underlying it, and not even the glaring helplessness of
this system in depicting the very important near regions of space has been a cause for
concern.  Sometimes the fetishization of this system went so far that it was declared the
only correct method of depiction (for one cannot argue with mathematics!), and hence
a mandatory hallmark of realistic painting.\footnote{%
  See, for example: N.\,N.~Volkov. \emph{Perception of Object and Drawing}.
  Moscow, 1950, pp.\,62, 181.  In his most recent monograph (N.\,N.~Volkov.
  \emph{Composition in Painting}. Moscow, 1977) the author has renounced such
  extreme assertions.

  In the article ``Perspective in the Picture'' (\emph{The Artist}, 1978, No.\,7)
  A.\,Zaytsev writes of linear perspective that it is a system ``permitting the artist to
  create on a two-dimensional plane the most convincing image of three-dimensional
  space'' (p.\,51).  The assertion about greatest convincingness has no basis and is
  moreover incorrect (this will be shown below, in Chapter~X\@), but such an assertion
  is typical of the position being discussed here.  No less expressive is the following
  reasoning: ``\ldots to see correctly in perspective is not easy\ldots the Old Russian
  artist\ldots natural perception led to a mass of incorrectnesses\ldots These were
  simply mistakes, because he drew as he saw'' (p.\,53).}

\srcpage{79}

It should be recalled that mathematics is merely a powerful instrument, and the results
of its use depend above all on what it is applied to.  The system of linear perspective
conveys only the first stage of visual perception, and its mathematical justification does
not go beyond this limit either.

In the age of the Renaissance, the great creators of the system of perspective had not
the slightest idea of the major contribution made by the work of the brain to the visual
perception system, and even had they known this, the level of development of
mathematics at the time would not have permitted it to be brought to bear on the
description of the workings of the ``eye~$+$~brain'' system.  To limit oneself to the
archaic conceptions of the fifteenth century today means dragging art history backwards.

The system of perceptual perspective can also be given rigorous mathematical
justification.  In doing so it proves necessary to use not school-level algebra and
geometry (as in the justification of the system of linear perspective), but integral
calculus, and in some cases even the geometry of Lobachevsky.  Consequently, the
system of perceptual perspective also has every right to be called a scientific
perspective.  Moreover~--- as should be the case in the exact sciences~--- the new,
perceptual scientific perspective does not supersede the old, but contains it within
itself as a special case (the depiction of distant regions of space).

The erroneous assessment, characteristic of the early twentieth century, of medieval art
as na\"ive and imperfect sometimes leads today to an unjustified apologetics of it and a
disparagement of the achievements of modern painting.  The harmfulness of both
extremes is easy to see.  It seems, however, that the fetishization of the system of
linear perspective~--- and not so much among art historians as among the broad public
of lovers of the visual arts~--- is especially harmful.
